\begin{ouverture}[college1]{4}{Solutions, ions et pH}{Partie 1 · Organisation et transformations de la matière}{Que signifie vraiment une solution acide ?}
\objectif{Relier pH, acidité et ions H+ ou OH−.}
\objectif{Distinguer dissolution, dilution et réaction.}
\objectif{Exploiter une banque de tests ioniques.}
\prerequis{Ion et charge.}
\prerequis{Soluté, solvant et solution.}
\end{ouverture}
\pcPage{Le cours}{Solutions, ions et pH}
\phantomsection\label{ess-4}
\essSection{1}{Des ions dans l’eau}
Une solution aqueuse contient de l’eau et des espèces dissoutes, molécules ou ions. Dans le modèle d’une solution de chlorure de sodium, les ions $\ce{Na+}$ et $\ce{Cl-}$ sont dispersés ; le sel n’a pas disparu. Une solution ionique peut conduire le courant grâce au déplacement des ions. Le courant dans un métal est lié à des électrons : les porteurs dépendent du milieu.\par
\pcFigure{figures/b04-solution.pdf}{Dissolution d’un solide ionique ; nombres d’ions dessinés choisis pour la lisibilité.}
\essSection{2}{Interpréter le pH}
Pour les solutions aqueuses usuelles vers 25 °C, pH inférieur à 7 : acide ; égal à 7 : neutre ; supérieur à 7 : basique. Les ions $\ce{H+}$ prédominent sur $\ce{OH-}$ dans une solution acide ; c’est l’inverse dans une solution basique. Les deux types d’ions peuvent être présents. Le pH est sans unité. Le papier pH donne une estimation ; le pH-mètre doit être réglé et utilisé selon son protocole.\par
\clearpage
\essSection{3}{Diluer ou faire réagir}
Diluer consiste à ajouter du solvant sans ajouter de soluté : une solution acide usuelle devient moins acide, son pH se rapproche de 7 ; une solution basique devient moins basique. Le pH n’est pas proportionnel à la quantité de solvant. Une réaction acide-base peut consommer les ions $\ce{H+}$ et $\ce{OH-}$ : $\ce{H+ + OH- -> H2O}$. La solution finale dépend des quantités introduites ; elle n’est pas automatiquement neutre.\par
\pcFigure{figures/b04-ph.pdf}{Échelle indicative pour des solutions aqueuses usuelles ; aucun goût ni contact direct.}
\essSection{4}{Utiliser des tests fournis}
Une banque de tests peut indiquer : avec nitrate d’argent, précipité blanc pour les ions chlorure ; avec soude, bleu pour $\ce{Cu^{2+}}$, vert pour $\ce{Fe^{2+}}$ ou brun pour $\ce{Fe^{3+}}$. On utilise des prélèvements séparés pour éviter de modifier tout le flacon. Le résultat s’interprète dans les conditions du test. Produits, protections et déchets sont gérés selon les consignes du professeur.\par
\begin{pcEncadre}{method}{Un exemple expliqué}Un liquide a pH 3. Il est acide ; dans le modèle, les ions H+ sont plus abondants que les ions OH−. Après dilution à l’eau, on mesure pH 4 : l’acidité diminue, mais le liquide reste acide. On ne déduit pas le nouveau pH par une moyenne des valeurs.\end{pcEncadre}
\pcPage{Expérimenter et retenir}{Solutions, ions et pH}
\phantomsection\label{activite-4}
\begin{pcEncadre}{method}{Investigation en classe}Le professeur fournit des solutions diluées autorisées et une banque de tests. Sur des prélèvements distincts, comparer papier pH et mesure instrumentale, puis identifier un ion par précipitation. Prévoir un témoin, noter la couleur observée et séparer observation, identification et limite du test.\end{pcEncadre}
\pcSousTitre{Ce que je dois savoir expliquer}\begin{questions}
\item Classer une solution avec son pH et citer les ions concernés.
\item Expliquer l’effet qualitatif d’une dilution.
\item Décrire un test avec réactif, observation et conclusion.
\end{questions}
\begin{pcEncadre}{attention}{Une erreur à éviter}Un pH de 4 n’est pas deux fois moins acide qu’un pH de 2. Au collège, on compare et on exploite les informations fournies sans traiter le pH comme une grandeur proportionnelle.\end{pcEncadre}
\pcRepere{cor-activite-4}{Pistes pour l’investigation}\pcGrille
\clearpage\markright{Je démarre}
\rubrique{Je démarre}
\par\noindent\begin{minipage}{\linewidth}
\begin{exercice}[1]{Classer}
\begin{questions}
\item Classer pH 2 et pH 11.
\item Que signifie pH 7 vers 25 °C ?
\end{questions}
\end{exercice}
\end{minipage}
\par\noindent\begin{minipage}{\linewidth}
\begin{exercice}[1]{L’unité}
\begin{questions}
\item Quelle est l’unité du pH ?
\item Quelle grandeur mesure-t-on avec une éprouvette ?
\end{questions}
\end{exercice}
\end{minipage}
\par\noindent\begin{minipage}{\linewidth}
\begin{exercice}[1]{Ions et acidité}
\begin{questions}
\item Quels ions prédominent dans une solution acide : $\ce{H+}$ ou $\ce{OH-}$ ?
\item Même question pour une solution basique.
\end{questions}
\end{exercice}
\end{minipage}
\par\noindent\begin{minipage}{\linewidth}
\begin{exercice}[1]{Du sel dissous}
\begin{questions}
\item Dans une solution de chlorure de sodium, citer les deux ions.
\item Pourquoi la filtration ordinaire ne les sépare-t-elle pas de l’eau ?
\end{questions}
\end{exercice}
\end{minipage}
\par\noindent\begin{minipage}{\linewidth}
\begin{exercice}[1]{Un test bleu}
\begin{questions}
\item La banque fournie associe soude + ions cuivre(II) à un précipité bleu. Que conclure si ce précipité apparaît ?
\item Faut-il verser le réactif dans le flacon d’origine ?
\end{questions}
\end{exercice}
\end{minipage}
\clearpage\markright{Je m’entraîne}
\rubrique{Je m’entraîne}
\par\noindent\begin{minipage}{\linewidth}
\begin{exercice}[2]{Une dilution}
\begin{questions}
\item Un acide usuel passe de pH 2 à pH 4 par ajout d’eau. Est-il encore acide ?
\item Son acidité a-t-elle augmenté ?
\end{questions}
\end{exercice}
\end{minipage}
\par\noindent\begin{minipage}{\linewidth}
\begin{exercice}[2]{Une base diluée}
\begin{questions}
\item Une base passe de pH 12 à pH 10. Comment évolue son caractère basique ?
\item Peut-on obtenir pH 5 par simple dilution idéale de cette base avec de l’eau neutre ?
\end{questions}
\end{exercice}
\end{minipage}
\par\noindent\begin{minipage}{\linewidth}
\begin{exercice}[2]{Trois inconnues}
\begin{questions}
\item A : pH 4 ; B : pH 7 ; C : pH 9. Classer.
\item Une solution contient des ions : peut-on déduire son pH sans les connaître ?
\end{questions}
\end{exercice}
\end{minipage}
\par\noindent\begin{minipage}{\linewidth}
\begin{exercice}[2]{Un résultat vert}
\begin{questions}
\item Le test à la soude donne un précipité vert, associé dans la banque à Fe2+. Nommer l’ion.
\item Combien d’électrons a-t-il perdus par rapport à l’atome ?
\end{questions}
\end{exercice}
\end{minipage}
\par\noindent\begin{minipage}{\linewidth}
\begin{exercice}[2]{Neutraliser}
\begin{questions}
\item Lire $\ce{H+ + OH- -> H2O}$ : quels ions sont consommés ?
\item Pourquoi un mélange acide-base ne finit-il pas toujours à pH 7 ?
\end{questions}
\end{exercice}
\end{minipage}
\clearpage\markright{Je raisonne}
\rubrique{Je raisonne}
\par\noindent\begin{minipage}{\linewidth}
\begin{exercice}[3]{Échantillons comparables}
\begin{questions}
\item Deux bandelettes indiquent environ 5 et environ 6. Peut-on annoncer une différence précise de 1,00 unité de pH ?
\item Proposer une amélioration.
\end{questions}
\end{exercice}
\end{minipage}
\par\noindent\begin{minipage}{\linewidth}
\begin{exercice}[3]{Conduction}
\begin{questions}
\item Comparer les porteurs du courant dans un fil métallique et dans une solution ionique.
\item La solution peut-elle rester globalement neutre tout en conduisant ?
\end{questions}
\end{exercice}
\end{minipage}
\par\noindent\begin{minipage}{\linewidth}
\begin{exercice}[3]{Sans moyenne}
\begin{questions}
\item Un élève annonce pH 7 en mélangeant des volumes égaux de solutions pH 2 et pH 12. Pourquoi ce calcul est-il insuffisant ?
\item Quelle donnée ou action manque ?
\end{questions}
\end{exercice}
\end{minipage}
\par\noindent\begin{minipage}{\linewidth}
\begin{exercice}[3]{L’eau et le climat}
\begin{questions}
\item Une eau passe de pH 8,2 à pH 8,0. Est-elle devenue acide ?
\item Peut-on parler d’acidification ?
\end{questions}
\end{exercice}
\end{minipage}
\par\noindent\begin{minipage}{\linewidth}
\begin{exercice}[3]{Concevoir une analyse}
\begin{questions}
\item Une solution peut contenir des ions chlorure et cuivre(II). Décrire les deux prélèvements et les tests nécessaires d’après la banque du cours.
\item Pourquoi un seul test ne suffit-il pas à conclure à la présence des deux ions ?
\end{questions}
\end{exercice}
\end{minipage}
\pcPage{Défi du chapitre}{Solutions, ions et pH}
\phantomsection\label{defi-ex-4}
\pcDefiTitre{4}{Le laboratoire enquête}
\begin{questions}
\item A : pH 3, test chlorure positif ; B : pH 10, test cuivre négatif. Classer A et B.
\item Un camarade veut mélanger A et B pour identifier tous les ions d’un coup. Pourquoi refuser ce protocole ?
\item Rédiger une conclusion de trois phrases limitée aux observations.
\end{questions}
\pcGrille
